% Options for packages loaded elsewhere
\PassOptionsToPackage{unicode}{hyperref}
\PassOptionsToPackage{hyphens}{url}
\PassOptionsToPackage{space}{xeCJK}
%
\documentclass[
]{ctexart}
\usepackage{amsmath,amssymb}
\usepackage{iftex}
\ifPDFTeX
  \usepackage[T1]{fontenc}
  \usepackage[utf8]{inputenc}
  \usepackage{textcomp} % provide euro and other symbols
\else % if luatex or xetex
  \usepackage{unicode-math} % this also loads fontspec
  \defaultfontfeatures{Scale=MatchLowercase}
  \defaultfontfeatures[\rmfamily]{Ligatures=TeX,Scale=1}
\fi
\usepackage{lmodern}
\ifPDFTeX\else
  % xetex/luatex font selection
  \ifXeTeX
    \usepackage{xeCJK}
    \setCJKmainfont[]{Songti SC}
          \fi
  \ifLuaTeX
    \usepackage[]{luatexja-fontspec}
    \setmainjfont[]{Songti SC}
  \fi
\fi
% Use upquote if available, for straight quotes in verbatim environments
\IfFileExists{upquote.sty}{\usepackage{upquote}}{}
\IfFileExists{microtype.sty}{% use microtype if available
  \usepackage[]{microtype}
  \UseMicrotypeSet[protrusion]{basicmath} % disable protrusion for tt fonts
}{}
\makeatletter
\@ifundefined{KOMAClassName}{% if non-KOMA class
  \IfFileExists{parskip.sty}{%
    \usepackage{parskip}
  }{% else
    \setlength{\parindent}{0pt}
    \setlength{\parskip}{6pt plus 2pt minus 1pt}}
}{% if KOMA class
  \KOMAoptions{parskip=half}}
\makeatother
\usepackage{xcolor}
\usepackage[margin=2cm]{geometry}
\setlength{\emergencystretch}{3em} % prevent overfull lines
\providecommand{\tightlist}{%
  \setlength{\itemsep}{0pt}\setlength{\parskip}{0pt}}
\setcounter{secnumdepth}{5}
\setcounter{section}{-1}
\ifLuaTeX
  \usepackage{selnolig}  % disable illegal ligatures
\fi
\IfFileExists{bookmark.sty}{\usepackage{bookmark}}{\usepackage{hyperref}}
\IfFileExists{xurl.sty}{\usepackage{xurl}}{} % add URL line breaks if available
\urlstyle{same}
\hypersetup{
  pdftitle={毛概期末划重点},
  pdfauthor={Rui},
  hidelinks,
  pdfcreator={LaTeX via pandoc}}

\title{毛概期末划重点}
\author{Rui}
\date{\texttt{r\ format(Sys.Date(),\ \textquotesingle{}\%B\ \%d,\ \%Y\textquotesingle{})}}

\begin{document}
\maketitle

\renewcommand*\contentsname{目录}
{
\setcounter{tocdepth}{3}
\tableofcontents
}
\hypertarget{ux5173ux4e8eux8003ux8bd5}{%
\section{关于考试}\label{ux5173ux4e8eux8003ux8bd5}}

\begin{itemize}
\tightlist
\item
  考试时间：2023.6.16（周五）18：30\textasciitilde20：20
\item
  填写学号和班级编号
\item
  考试题型

  \begin{itemize}
  \tightlist
  \item
    四道简答或辨析题：每道 6 分
  \item
    两道材料分析题：每道 13 分
  \end{itemize}
\item
  成绩构成：

  \begin{itemize}
  \tightlist
  \item
    期末考试：50\%
  \item
    平时成绩：50\%（考勤10\%+课堂展示讨论10\%+课程论文30\%）
  \end{itemize}
\item
  考察范围：以教材和授课PPT为主，同时了解党和国家最新的理论动向以及时政热点
\item
  答题要求：分段落、分要点，逻辑条理清晰，以马克思主义中国化时代化相关理论为基础阐明自己的观点，言之有理。
\end{itemize}

\hypertarget{ux515aux7684ux4e8cux5341ux5927}{%
\section{党的二十大}\label{ux515aux7684ux4e8cux5341ux5927}}

\begin{itemize}
\tightlist
\item
  三个务必

  \begin{itemize}
  \tightlist
  \item
    全党同志务必不忘初心、牢记使命
  \item
    务必谦虚谨慎、艰苦奋斗
  \item
    务必敢于斗争、善于斗争
  \end{itemize}
\item
  三件大事

  \begin{itemize}
  \tightlist
  \item
    一是迎来中国共产党成立一百周年
  \item
    二是中国特色社会主义进入新时代
  \item
    三是完成脱贫攻坚、全面建成小康社会的历史任务，实现第一个百年奋斗目标
  \end{itemize}
\end{itemize}

\hypertarget{ux4eceux4e94ux5e74ux89c4ux5212ux52302035ux5e74ux603bux4f53ux76eeux6807ux4e2dux56fdux5f0fux73b0ux4ee3ux5316}{%
\section{从五年规划到2035年总体目标------中国式现代化}\label{ux4eceux4e94ux5e74ux89c4ux5212ux52302035ux5e74ux603bux4f53ux76eeux6807ux4e2dux56fdux5f0fux73b0ux4ee3ux5316}}

``中国共产党的中心任务就是团结带领全国各族人民全面建成社会主义现代化强国、实现第二个百年奋斗目标，以中国式现代化全面推进中华民族伟大复兴。''

\begin{itemize}
\tightlist
\item
  中国式现代化的五个特征

  \begin{itemize}
  \tightlist
  \item
    人口规模巨大
  \item
    全体人民共同富裕
  \item
    物质文明和精神文明相协调
  \item
    人与自然和谐共生
  \item
    走和平发展道路
  \end{itemize}
\item
  老三步走：党的十三大明确提出``三步走''的发展战略

  \begin{itemize}
  \tightlist
  \item
    第一步从1981年到1990年国民生产总值翻一番，解决人民的温饱问题
  \item
    第二步从1991年到20世纪末，国民生产总值再翻一番，人民生活达到小康水平
  \item
    第三步到21世纪中叶，人均国民生产总值达到中等发达国家水平，人民生活比较富裕，基本实现现代化
  \end{itemize}
\item
  新三步走

  \begin{itemize}
  \tightlist
  \item
    第一步，2001-2010年时期，实现GNP翻一番，使人民的小康生活更加宽裕，形成比较完善的社会主义市场经济体制（十五大，1997）；使经济总量、综合国力和人民生活水平再上一个大台阶，为后十年的更大发展打好基础（十六大，2002）
  \item
    第二步，2010-2020年时期，GDP翻两番，人均GDP达到3000美元，全面建设小康社会，综合国力和争力明显增强
    （十六大，2002）
  \item
    第三步，2020-2050年时期，基本实现现代化，把我国建成富强民主文明的社会主义国家
    （十六大，2002）
  \end{itemize}
\item
  ``两个一百年''的奋斗目标

  \begin{itemize}
  \tightlist
  \item
    到建党100年时，建成惠及十几亿人口的更高水平的小康社会
  \item
    到新中国成立100年时，基本实现现代化，把我国建成社会主义现代化国家
  \end{itemize}
\item
  两步走''战略

  \begin{itemize}
  \tightlist
  \item
    第一阶段：从二O二O年到二O三五年，在全面建成小康社会的基础上，再奋斗十五年，基本实现社会主义现代化
  \item
    第二阶段：从二O三五年到本世纪中叶，在基本实现现代化的基础上，再奋斗十五年，把我国建成富强民主文明和谐美丽的社会主义现代化强国
  \end{itemize}
\item
  2023年总体目标的具体内容

  \begin{enumerate}
  \def\labelenumi{\arabic{enumi}.}
  \tightlist
  \item
    经济实力、科技实力、综合国力大幅跃升，人均国内生产总值迈上新的大台阶，达到中等发达国家水平
  \item
    实现高水平科技自立自强，进入创新型国家前列；
  \item
    建成现代化经济体系，形成新发展格局，基本实现新型工业化、信息化、城镇化、农业现代化
  \item
    基本实现国家治理体系和治理能力现代化，全过程人民民主制度更加健全，基本建成法治国家、法治政府、法治社会
  \item
    建成教育强国、科技强国、人才强国、文化强国、体育强国、健康中国，国家文化软实力显著增强
  \item
    人民生活更加幸福美好，居民人均可支配收入再上新台阶，中等收入群体比重明显提高，基本服务实现均等化，农村基本具备现代生活条件，社会保持长期稳定，人的全面发展、全体人司富裕取得更为明显的实质性进展
  \item
    广泛形成绿色生产生活方式，碳排放达峰后稳中有降，生态环境根本好转，美丽中国目标基本实现
    8.国家安全体系和能力全面加强，基本实现国防和军队现代化
  \end{enumerate}
\end{itemize}

\hypertarget{ux9ad8ux8d28ux91cfux53d1ux5c55}{%
\section{高质量发展}\label{ux9ad8ux8d28ux91cfux53d1ux5c55}}

（结合社会主义市场经济）

\begin{itemize}
\tightlist
\item
  全面建设社会主义现代化国家的首要任务，要坚持以推动高质量发展为主题。
\item
  推动经济社会发展绿色化、低碳化是实现高质量发展的关键环节。
\item
  坚持把发展经济的着力点放在实体经济上。
\end{itemize}

\hypertarget{ux793eux4f1aux4e3bux4e49ux5e02ux573aux7ecfux6d4e}{%
\section{社会主义市场经济}\label{ux793eux4f1aux4e3bux4e49ux5e02ux573aux7ecfux6d4e}}

第七讲：邓小平理论的基本问题和主要内容

\begin{itemize}
\tightlist
\item
  党的十二届三中全会通过《中共中央关于经济体制改革的决定》提出了社会主义经济是``公有制基础上的有计划的商品经济''的论断
\item
  南方谈话中，邓小平明确提出：``计划经济不等于社会主义，资本主义计划也有计划；市场经济不等于资本主义，社会主义也有市场''
\item
  党的十四大明确提出，我国经济体制改革的目标是建立社会主义市场经济体制
\end{itemize}

在社会主义市场经济条件下规范和引导资本发展，既是一个重大经济问题也是一个重大政治问题，既是一个重大实践问题，也是一个重大理论问题关系坚持社会主义基本经济制度，关系改革开放基本国策，关系高质量发展和共同富裕，关系国家安全和社会稳定。\textbf{资本作为重要生产要素，是市场配置资源的工具，是发展经济的方式和手段，社会主义国家也可以利用各类资本推动经济社会发展。}现阶段，我国存在国有资本、集体资本、民营资本、外国资本、混合资本等各种形态资本，并呈现出规模显著增加、主体更加多元、运行速度加快、国际资本大量进入等明显特征。

\hypertarget{ux5bb6ux5eadux8054ux4ea7ux627fux5305ux8d23ux4efbux5236ux4e61ux6751ux632fux5174}{%
\section{家庭联产承包责任制------乡村振兴}\label{ux5bb6ux5eadux8054ux4ea7ux627fux5305ux8d23ux4efbux5236ux4e61ux6751ux632fux5174}}

\begin{itemize}
\tightlist
\item
  历史沿革

  \begin{itemize}
  \tightlist
  \item
    第一个中央
    ``一号文件``：1982年1月1日，中共中央批转1981年12月的《全国农村工作会议纪要》，主要内容就是肯定多种形式的责任制，特别是包干到户、包产到户，正式承认包产到户合法性。
  \item
    第二个中央
    ``一号文件''：1983年1月2日，中共中央印发《当前农村经济政策的若干问题》文件从理论上说明了家庭联产承包责任制``是在党的领导下中国农民的伟大创造，是马克思主义农业合作化理论在我国实践中的新发展''；提出了``两个转化''，即促进农业从自给半自给经济向较大规模的商品生产转化，从传统农业向现代农业转化，从而放活农村工商业。
  \item
    第三个中央``一号文件''：1984年1月1日《关于1984年农村工作的通知》中提出，延长土地承包期，土地承包期一般应在十五年以上\ldots..允许有偿转让土地使用权；鼓励农民向各种企业投资入股；继续减少统派购的品种和数量；允许务工、经商、办服务业的农民自理口粮到集镇落户。文件确立农村工作的重点是：在稳定和完善生产责任制的基础上，提高生产力水平，疏理流通渠道，发展农村商品生产。
  \end{itemize}
\item
  总目标：实现农业农村现代化
\item
  总方针：坚持农业农村优先发展
\item
  总要求：产业兴旺、生态宜居、乡风文明、治理有效、生活富裕
\item
  实施乡村振兴战略是建设现代化经济体系的重要基础
\item
  乡村振兴是关系全面建设社会主义现代化国家的全局性、历史性任务，是新时代``三农''工作总抓手
\item
  中共中央办公厅、国务院办公厅印发《乡村建设行动实施方案》：乡村建设是实施乡村振兴战略的重要任务，也是国家现代化建设的重要内容。坚持农业农村优先发展，把乡村建设摆在社会主义现代化建设的重要位置,顺应农民群众对美好生活的向往，以普惠性、基础性、兜底性民生建设为重点。
\end{itemize}

\hypertarget{ux5148ux5bccux5e26ux540eux5bccux5171ux540cux5bccux88d5}{%
\section{先富带后富------共同富裕}\label{ux5148ux5bccux5e26ux540eux5bccux5171ux540cux5bccux88d5}}

\begin{itemize}
\tightlist
\item
  《求是》2021年第20期发表的习总书记重要文章《扎实推动共同富裕》：共同富裕是社会主义的本质要求，是中国式现代化的重要特征。
\item
  共同富裕是全体人民共同富裕，是人民群众物质生活和精神生活都富裕；共同富裕是一个长远目标；不同人群不仅实现富裕的程度有高有低，时间上也会有先有后，不同地区富裕程度还会存在一定差异。这是一个在动态中向前发展的过程，要持续推动，不断取得成效。
\item
  共同富裕不是少数人的富裕，也不是整齐划一的平均主义；不是所有人都同时富裕，也不是所有地区同时达到一个富裕水准；不可能一蹴而就；不可能齐头并进。
\end{itemize}

\hypertarget{ux7edfux4e00ux6218ux7ebf}{%
\section{统一战线}\label{ux7edfux4e00ux6218ux7ebf}}

第二讲：毛泽东思想及其历史地位

遵义会议以后，毛泽东系统地总结了中国革命的经验，分析和批判了教条主义的错误，并及时吸取抗日战争的新经验，形成了比较系统的哲学、军事、统一战线和党的建设思想，系统地阐述了中国新民主主义革命的基本理论、路线和纲领，精辟地论证了党在民主革命时期的政策和策略。

\begin{itemize}
\tightlist
\item
  国民革命时期：国共统一战线。由广大工人、农民、城市小资产阶级参加的反帝反封建的国民革命联合战线
\item
  土地革命时期：工农民主统一战线。包括两个联盟：工、农、知识分子和其他劳动者的联盟；工、农、知识分子和非劳动者的联盟
\item
  抗日战争时期：抗日民族统一战线(第二次国共合作)。包括了工人阶级、农民阶级、小资产阶级和民族资产阶级，蒋介石为首的国民党亲英美派
\item
  解放战争时期：人民民主统一战线。包括各民族、各民主阶级、各民主党派、各人民团体、广大华侨、各界民主人士及其他爱国分子和国民党统治集团中的一部分地方实力派
\item
  改革开放时期：十一届三中全会后，提出爱国统一战线。它是工人阶级领导的、工农联盟为基础的全体社会主义劳动者、社会主义建设者、拥护社会主义的爱国者和拥护祖国统一的爱国者的联盟，是最广泛的爱国统一战线
\item
  新时代统一战线的定位：

  \begin{enumerate}
  \def\labelenumi{\arabic{enumi}.}
  \tightlist
  \item
    以服务中华民族伟大复兴为突出功能
  \item
    以实现中华民族伟大复兴为最大价值公约数
  \item
    以展现中华民族伟大复兴的中国方案为外部效应
  \end{enumerate}
\item
  新时代统一战线的发展：

  \begin{enumerate}
  \def\labelenumi{\arabic{enumi}.}
  \tightlist
  \item
    形成国内与国际、网上与网下高度统筹的统一战线空间布局
  \item
    形成策略与战略、法宝与目的高度一致的统一战线理论认识
  \item
    形成利益与认同、政策与法治高度协同的统一战线工作方法
  \item
    形成衔接紧密、权威高效的统一战线运转体制
  \item
    形成紧贴时代、着眼未来的统一战线发展视野
  \end{enumerate}
\end{itemize}

\hypertarget{ux4e00ux56fdux4e24ux5236ux6e2fux6fb3ux53f0}{%
\section{一国两制------港澳台}\label{ux4e00ux56fdux4e24ux5236ux6e2fux6fb3ux53f0}}

第七讲：邓小平理论的主要内容

\hypertarget{ux9999ux6e2fux6fb3ux95e8}{%
\subsection{香港澳门}\label{ux9999ux6e2fux6fb3ux95e8}}

\begin{itemize}
\tightlist
\item
  ``一国两制''的提出和意义：1982年1月11日，邓小平首次提出``一种国家两种制度''的概念。``一国两制''伟大构想，把和平共处的原则用之于解决一个国家的统一问题，是对马克思主义国家学说的创造性发展。``一国两制''作为基本国策，成为党领导人民实现祖国和平统一的一项重要制度和中国特色社会主义的一个伟大创举，在实现和维护国家统一、维护国家安全、实现中华民族伟大复兴、推进中国特色社会主义事业等方面作出了重要贡献。
\item
  维护国家主权、安全、发展利益是``一国两制''方针的最高原则，把香港、澳门特别行政区管治权牢牢掌握在爱国者手中，是保证香港、澳门长治久安的必然要求，``一国两制''是特别行政区宪制秩序的根本遵循，中央对特别行政区拥有全面管治权是特别行政区高度自治权的源头。
\item
  ``一国两制''的伟大构想的提出是从解决台湾问题开始的：台湾问题是国内战争遗留下来的问题，属于中国的内政，不容许外国干涉。
\item
  ``一国两制''的伟大构想在实践中首先运用于解决香港问题、澳门问题：港澳问题是殖民主义的遗留问题，分属中英、中葡之间的问题
\item
  ``和平统一，一国两制''

  \begin{itemize}
  \tightlist
  \item
    核心坚持一个中国；两制并存一一在祖国统一的前提下，国家的主体部分实行社会主义制度，同时在台湾、香港、澳门保持原有的社会制度和生活方式长期不变
  \item
    高度自治一-祖国完全统一后，台湾、香港、澳门作为特别行政区，享有不同于中国其他省、市、自
    治区的高度自治权，台湾、香港、澳门同胞各种合法权益将得到切实尊重和维护
  \end{itemize}
\item
  ``一国两制''方针的最高原则：维护国家主权、安全、发展利益。
\item
  解决台湾问题、实现祖国完全统一，是党矢志不渝的历史任务，是全体中华儿女的共同愿望,是实现中华民族伟大复兴的必然要求。
\item
  证明中国从法律和事实上收复了台湾的一系列具有国际法律效力的文件：《开罗宣言》《波茨坦公告》、《日本投降条款》。
\end{itemize}

\hypertarget{ux53f0ux6e7e}{%
\subsection{台湾}\label{ux53f0ux6e7e}}

\begin{itemize}
\tightlist
\item
  美国国会众议长佩洛西不顾中方强烈反对和严正交涉，窜访中国台湾地区严重违反一个中国原则和中美三个联合公报规定，严重冲击中美关系政治基础，严重侵犯中国主权和领土完整，严重破坏台海和平稳定，向``台独''势力发出严重错误信号。``台独''分裂行径是祖国统一的最大障碍，是民族复兴的严重隐患。
\item
  《台湾问题与新时代中国统一事业白皮书》指出，中华民族迎来了从站起来富起来到强起来的伟大飞跃，实现中华民族伟大复兴进入了不可逆转的历史进程。这是中国统一大业新的历史方位。
\item
  ``和平统一，一国两制''是我们解决台湾问题的基本方针，也是实现祖国统一的最佳方式，既充分考虑台湾现实情况，又有利于统一后台湾长治久安。
\end{itemize}

\hypertarget{ux79d1ux5b66ux6280ux672fux662fux7b2cux4e00ux751fux4ea7ux529b}{%
\section{科学技术是第一生产力}\label{ux79d1ux5b66ux6280ux672fux662fux7b2cux4e00ux751fux4ea7ux529b}}

\begin{itemize}
\tightlist
\item
  1978年3月全国科学大会展开，迎来了科学界的春天。邓小平关于``科学技术是第一生产力''的重大论断，为我们国家和人民掌握科学技术，赶上全球的发展，赶上世界的步伐，指明了努力的方向。
\item
  ``863计划''

  \begin{itemize}
  \tightlist
  \item
    1986年3月3日，王大、王昌、杨嘉握、陈芳允四位科学家向国家提出要跟踪世界先进水平，发展中国高技术的建议。
  \item
    《高技术发展研究纲要》
  \end{itemize}
\item
  事例

  \begin{itemize}
  \tightlist
  \item
    C919大型客机取得型号合格证，标志着我国具备按照国际通行适航标准研制大型客机的能力。
  \item
    2020年，北斗三号系统正式建成开通，面向全球提供卫星导航服务，标志着北斗系统进入全球化发展新阶段，标志着北斗系统``三步走''发展战略圆满完成
  \item
    中国载人航天工程三步走立项实施30周年：第三步是计划在2022年底之前建成长期载人的大型空间站``天宫''，开展大规模、长期有人照料的空间应用。2022年为空间站在轨建造阶段，2022年7月24日发射了``问天''实验舱，2022年10月31日发射了``梦天实验舱，两个实验舱与``天和''核心舱形成``T''字基本构型，建成国家太空实验室此后为``天宫''空间站运行阶段。
  \end{itemize}
\end{itemize}

\hypertarget{ux4e09ux4e2aux4ee3ux8868ux91cdux8981ux601dux60f3}{%
\section{``三个代表''重要思想}\label{ux4e09ux4e2aux4ee3ux8868ux91cdux8981ux601dux60f3}}

第八讲：``三个代表''重要思想及其历史地位；第九讲：``三个代表''重要思想的主要内容

\begin{itemize}
\item
  江泽民同志是我党我军我国各族人民公认的享有崇高威望的卓越领导人，伟大的马克思主义者，伟大的无产阶级革命家、政治家、军事家、外交家、久经考验的共产主义战士，中国特色社会主义伟大事业的杰出领导者，党的第三代中央领导集体的核心，``三个代表''重要思想的主要创立者。
\item
  江泽民同志的成就：江泽民同志带领党的中央领导集体，紧紧依靠全党全军全国各族人民，旗帜鲜明坚持四项基本原则，维护国家独立、尊严、安全、稳定，毫不动摇坚持经济建设这个中心，坚持改革开放，捍卫了中国特色社会主义伟大事业，打开了我国改革开放和社会主义现代化建设新局面，确立了社会主义市场经济体制的改革目标和基本框架，确立了社会主义初级阶段公有制为主体、多种所有制经济共同发展的基本经济制度和按劳分配为主体、多种分配方式并存的分配制度，制定和实施了促进改革发展稳定的一系列方针政策和重大战略，锐意推进经济体制改革、政治体制改革、文化体制改革和其他各方面改革，开创全面改革新局面，实施依法治国基本方略，坚持``和平统一、一国两制''的方针，实现香港、澳门顺利回归,坚持独立自主的和平外交政策，打开外交工作崭新局面，推进党的建设新的伟大工程，推动社会主义物质文明、政治文明、精神文明建设取得举世瞩目的新进展，成功把中国特色社会主义推向二十一世纪。
\item
  ``三个代表''重要思想的核心观点：中国共产党必须始终代表中国先进生产力的发展要求，代表中国先进文化的前进方向，代表中国最广大人民的根本利益。这是我们党的立党之本、执政之基、力量之源。
\item
  ``三个代表''重要思想最鲜明的特点和最突出的贡献：在于用一系列紧密联系相互贯通的新思想新观点新论断，进一步回答了什么是社会主义、怎样建设社会主义的问题，创造性回答了建设什么样的党、怎样建设党的问题，是对马克思列宁主义、毛泽东思想、邓小平理论的继承和发展，深化了我们对新的时代条件下推进中国特色社会主义事业、加强党的建设的规律的认识。
\end{itemize}

\end{document}
